\documentclass[10pt,a4paper]{amsart}
\usepackage{amsmath,amssymb,amsthm,graphicx}
% Retained preamble for the Stacks Project chapter divisors.tex (Divisors).
% The chapter file pinned-original-source.tex carries only \input{preamble} on its first line, so no
% preamble text lives inside the chapter itself.  The Stacks Project uses one single shared preamble for
% every chapter; the snapshot stages it verbatim in sources/stacks-chow/preamble.tex.  The block below is
% copied mechanically, line for line, from that file's lines 14-18 (the Xy-pic setup) and lines 154-217
% (the native theorem declarations together with the mathematical macros).  Nothing was added, removed or
% redefined.  The wrapper class and the amsmath/amssymb/amsthm/graphicx base packages are supplied by the
% assembler, so the class/\externaldocument part of the shared preamble is not repeated here.
% For commutative diagrams we use Xy-pic
\usepackage[all]{xy}

% We use 2cell for 2-commutative diagrams.
\xyoption{2cell}
% Theorem environments.
%
\theoremstyle{plain}
\newtheorem{theorem}[subsection]{Theorem}
\newtheorem{proposition}[subsection]{Proposition}
\newtheorem{lemma}[subsection]{Lemma}

\theoremstyle{definition}
\newtheorem{definition}[subsection]{Definition}
\newtheorem{example}[subsection]{Example}
\newtheorem{exercise}[subsection]{Exercise}
\newtheorem{situation}[subsection]{Situation}

\theoremstyle{remark}
\newtheorem{remark}[subsection]{Remark}
\newtheorem{remarks}[subsection]{Remarks}

\numberwithin{equation}{subsection}

% Macros
%
\def\lim{\mathop{\mathrm{lim}}\nolimits}
\def\colim{\mathop{\mathrm{colim}}\nolimits}
\def\Spec{\mathop{\mathrm{Spec}}}
\def\Hom{\mathop{\mathrm{Hom}}\nolimits}
\def\Ext{\mathop{\mathrm{Ext}}\nolimits}
\def\SheafHom{\mathop{\mathcal{H}\!\mathit{om}}\nolimits}
\def\SheafExt{\mathop{\mathcal{E}\!\mathit{xt}}\nolimits}
\def\Sch{\mathit{Sch}}
\def\Mor{\mathop{\mathrm{Mor}}\nolimits}
\def\Ob{\mathop{\mathrm{Ob}}\nolimits}
\def\Sh{\mathop{\mathit{Sh}}\nolimits}
\def\NL{\mathop{N\!L}\nolimits}
\def\CH{\mathop{\mathrm{CH}}\nolimits}
\def\proetale{{pro\text{-}\acute{e}tale}}
\def\etale{{\acute{e}tale}}
\def\QCoh{\mathit{QCoh}}
\def\Ker{\mathop{\mathrm{Ker}}}
\def\Im{\mathop{\mathrm{Im}}}
\def\Coker{\mathop{\mathrm{Coker}}}
\def\Coim{\mathop{\mathrm{Coim}}}

% Boxtimes
%
\DeclareMathSymbol{\boxtimes}{\mathbin}{AMSa}{"02}

%
% Macros for moduli stacks/spaces
%
\def\QCohstack{\mathcal{QC}\!\mathit{oh}}
\def\Cohstack{\mathcal{C}\!\mathit{oh}}
\def\Spacesstack{\mathcal{S}\!\mathit{paces}}
\def\Quotfunctor{\mathrm{Quot}}
\def\Hilbfunctor{\mathrm{Hilb}}
\def\Curvesstack{\mathcal{C}\!\mathit{urves}}
\def\Polarizedstack{\mathcal{P}\!\mathit{olarized}}
\def\Complexesstack{\mathcal{C}\!\mathit{omplexes}}
% \Pic is the operator that assigns to X its picard group, usage \Pic(X)
% \Picardstack_{X/B} denotes the Picard stack of X over B
% \Picardfunctor_{X/B} denotes the Picard functor of X over B
\def\Pic{\mathop{\mathrm{Pic}}\nolimits}
\def\Picardstack{\mathcal{P}\!\mathit{ic}}
\def\Picardfunctor{\mathrm{Pic}}
\def\Deformationcategory{\mathcal{D}\!\mathit{ef}}

% --- appended 2026-10-12 (Chief): the project own definitions of the reference/slogan/history
% display environments, copied mechanically from the shared Stacks preamble lines 8-12
% (\usepackage{verbatim} plus the three \newenvironment lines, which the chapter preamble keeps
% verbatim from the project preamble). Nothing added, removed or redefined: these environments
% carry no mathematical content and the project itself typesets them as comments. Without these
% three lines a counted statement whose own text uses them fails with Environment slogan undefined.
\usepackage{verbatim}
\newenvironment{reference}{\comment}{\endcomment}
\newenvironment{slogan}{\comment}{\endcomment}
\newenvironment{history}{\comment}{\endcomment}
\usepackage{cleveref}
\crefname{originallemma}{original lemma}{original lemmas}
\Crefname{originallemma}{Original lemma}{Original lemmas}
\makeatletter
\newlabel{algebra-lemma-characterize-dvr}{{00PD}{source}{}{sourcepointer-1}{}}
\newlabel{algebra-lemma-characterize-dvr@cref}{{[originallemma][0][]00PD}{[1][0][]source}{}{}{}}
\makeatother
\crefname{originallemma}{original lemma}{original lemmas}
\Crefname{originallemma}{Original lemma}{Original lemmas}
\makeatletter
\newlabel{algebra-lemma-finite-projective}{{00NX}{source}{}{sourcepointer-2}{}}
\newlabel{algebra-lemma-finite-projective@cref}{{[originallemma][0][]00NX}{[1][0][]source}{}{}{}}
\makeatother
\crefname{originallemma}{original lemma}{original lemmas}
\Crefname{originallemma}{Original lemma}{Original lemmas}
\makeatletter
\newlabel{algebra-lemma-regular-mcm-free}{{00NT}{source}{}{sourcepointer-3}{}}
\newlabel{algebra-lemma-regular-mcm-free@cref}{{[originallemma][0][]00NT}{[1][0][]source}{}{}{}}
\makeatother
\crefname{originallemma}{original lemma}{original lemmas}
\Crefname{originallemma}{Original lemma}{Original lemmas}
\makeatletter
\newlabel{more-algebra-lemma-dedekind-torsion-free-flat}{{0AUW}{source}{}{sourcepointer-4}{}}
\newlabel{more-algebra-lemma-dedekind-torsion-free-flat@cref}{{[originallemma][0][]0AUW}{[1][0][]source}{}{}{}}
\makeatother
\crefname{originallemma}{original lemma}{original lemmas}
\Crefname{originallemma}{Original lemma}{Original lemmas}
\makeatletter
\newlabel{more-algebra-lemma-reflexive-over-normal}{{0AY6}{source}{}{sourcepointer-5}{}}
\newlabel{more-algebra-lemma-reflexive-over-normal@cref}{{[originallemma][0][]0AY6}{[1][0][]source}{}{}{}}
\makeatother
\title{On a regular scheme of dimension at most two a coherent module is reflexive exactly when it is finite locally free}
\author{The Stacks Project Authors}
\date{stacks-project:divisors.tex:local-snapshot-20261012; GNU FDL 1.2 or later}
\begin{document}
\maketitle
\noindent\textit{Editorial note.} All delivered mathematical text is the author's own, copied from the pinned Stacks Project chapter file pinned-original-source.tex; no statement or proof step was rewritten, reordered or removed.\par\medskip
\noindent\textit{Editorial note.} Following the multi-claim convention for this source, source\_url is the source-level URL https://stacks.math.columbia.edu/ and the source label lemma-reflexive-over-regular-dim-2 is carried in claim\_key and here rather than in a per-claim tag URL.\par\medskip
\noindent\textit{Editorial note.} The preamble is the chapter's shared Stacks Project preamble, retained as sources/stacks-divisors/retained-preamble.tex: a line-for-line excerpt of the shared Stacks preamble staged with this snapshot. No macro was invented or redefined, and this card adds no preamble text of its own.\par\medskip
\noindent\textit{Editorial note.} Statement span 2336-2345 and proof span 2347-2368 come from the halfA pool card operations/pool-cards/stacks-divisors-halfA.json and were re-read in the pinned chapter file before this card was written.\par\medskip
\noindent\textit{Editorial note.} The proof splits into the three possible values of the dimension of the local ring and refers to the cited statements by their hypotheses; the author's wording is preserved.\par\medskip
\noindent\textit{Editorial note.} 5 cross-references occur inside the counted spans; they are listed in external\_references as navigation pointers only.\par\medskip
\noindent\textit{Allowed standard prerequisites (editorial).}\begin{itemize}
\item A finite locally free coherent module is reflexive. Used at proof lines 2347-2348. Gives the easy direction.
\item A coherent module whose stalks are free is finite locally free, cited from the algebra chapter. Cited at proof lines 2350-2352. Reduces the converse to freeness of the stalk at each point.
\item Local rings of dimension zero are fields and local rings of dimension one that are regular are discrete valuation rings, with torsion free finite modules over them free, cited from the algebra and more-algebra chapters. Cited at proof lines 2353-2359. Disposes of the points of dimension at most one.
\item A reflexive module over a normal local ring has depth at least two, and a maximal Cohen-Macaulay module over a regular local ring is free, cited from the more-algebra and algebra chapters. Cited at proof lines 2360-2367. Disposes of the points of dimension two.
\end{itemize}\medskip
\ifdefined\hypertarget\hypertarget{sourcepointer-1}{}\fi
\noindent\textit{Source reference.} Original lemma 00PD: focused-seeds/AGP/sources/stacks-snapshot-20261011/tags.txt, source lines 927--927. Same-chapter reference used in the retained author text at source line(s) 1819, 2358, 3021, 7023. It is supplied only as a navigation pointer so that the retained text compiles; the counted item delivers the author's own statement and the author's complete proof, and no part of the referenced item is copied in or substituted for the counted argument.\par\medskip
\ifdefined\hypertarget\hypertarget{sourcepointer-2}{}\fi
\noindent\textit{Source reference.} Original lemma 00NX: focused-seeds/AGP/sources/stacks-snapshot-20261011/tags.txt, source lines 875--875. Same-chapter reference used in the retained author text at source line(s) 1813, 2352. It is supplied only as a navigation pointer so that the retained text compiles; the counted item delivers the author's own statement and the author's complete proof, and no part of the referenced item is copied in or substituted for the counted argument.\par\medskip
\ifdefined\hypertarget\hypertarget{sourcepointer-3}{}\fi
\noindent\textit{Source reference.} Original lemma 00NT: focused-seeds/AGP/sources/stacks-snapshot-20261011/tags.txt, source lines 871--871. Same-chapter reference used in the retained author text at source line(s) 2367. It is supplied only as a navigation pointer so that the retained text compiles; the counted item delivers the author's own statement and the author's complete proof, and no part of the referenced item is copied in or substituted for the counted argument.\par\medskip
\ifdefined\hypertarget\hypertarget{sourcepointer-4}{}\fi
\noindent\textit{Source reference.} Original lemma 0AUW: focused-seeds/AGP/sources/stacks-snapshot-20261011/tags.txt, source lines 13664--13664. Same-chapter reference used in the retained author text at source line(s) 1822, 2293, 2361. It is supplied only as a navigation pointer so that the retained text compiles; the counted item delivers the author's own statement and the author's complete proof, and no part of the referenced item is copied in or substituted for the counted argument.\par\medskip
\ifdefined\hypertarget\hypertarget{sourcepointer-5}{}\fi
\noindent\textit{Source reference.} Original lemma 0AY6: focused-seeds/AGP/sources/stacks-snapshot-20261011/tags.txt, source lines 13779--13779. Same-chapter reference used in the retained author text at source line(s) 2275, 2364. It is supplied only as a navigation pointer so that the retained text compiles; the counted item delivers the author's own statement and the author's complete proof, and no part of the referenced item is copied in or substituted for the counted argument.\par\medskip
\begin{lemma}
\label{lemma-reflexive-over-regular-dim-2}
Let $X$ be a regular scheme of dimension $\leq 2$.
Let $\mathcal{F}$ be a coherent $\mathcal{O}_X$-module.
The following are equivalent
\begin{enumerate}
\item $\mathcal{F}$ is reflexive,
\item $\mathcal{F}$ is finite locally free.
\end{enumerate}
\end{lemma}

\begin{proof}
It is clear that a finite locally free module is reflexive.
For the converse, we will show that if $\mathcal{F}$ is
reflexive, then $\mathcal{F}_x$ is a free $\mathcal{O}_{X, x}$-module
for all $x \in X$. This is enough by
Algebra, Lemma \ref{algebra-lemma-finite-projective}
and the fact that $\mathcal{F}$ is coherent.
If $\dim(\mathcal{O}_{X, x}) = 0$, then
$\mathcal{O}_{X, x}$ is a field and the statement is clear.
If $\dim(\mathcal{O}_{X, x}) = 1$, then $\mathcal{O}_{X, x}$
is a discrete valuation ring
(Algebra, Lemma \ref{algebra-lemma-characterize-dvr})
and $\mathcal{F}_x$ is torsion free.
Hence $\mathcal{F}_x$ is free by More on Algebra, Lemma
\ref{more-algebra-lemma-dedekind-torsion-free-flat}.
If $\dim(\mathcal{O}_{X, x}) = 2$, then $\mathcal{O}_{X, x}$
is a regular local ring of dimension $2$. By
More on Algebra, Lemma \ref{more-algebra-lemma-reflexive-over-normal}
we see that $\mathcal{F}_x$ has depth $\geq 2$.
Hence $\mathcal{F}_x$ is free by
Algebra, Lemma \ref{algebra-lemma-regular-mcm-free}.
\end{proof}

\end{document}
